% ============================================================================
% WORKBOOK — Section 6.1: Understanding and Using Scale Drawings
% CCSS 7.G.A.1
% Practice-focused reinforcement for the study guide topic
% ============================================================================
\section{Understanding and Using Scale Drawings}
\topicTitle{Understanding and Using Scale Drawings}

\begin{quickReview}{Scale Drawings}
A \textbf{scale drawing} is a proportional copy of a real object where every length is multiplied by the same \textbf{scale factor}.

\begin{itemize}[leftmargin=*]
    \item \textbf{Actual length} $=$ drawing length $\times$ scale factor.
    \item If the scale is ``$1$ cm $= 5$ m,'' the scale factor is $5$. A $3$ cm line on the drawing represents $3 \times 5 = 15$ m.
    \item \textbf{Area scales by $k^2$:} if lengths scale by $k$, then areas scale by $k^2$.
\end{itemize}

\medskip
\textbf{Example:} A blueprint uses $1$ in $= 6$ ft. A wall is $4$ in on the blueprint. Actual length: $4 \times 6 = 24$ ft.
\end{quickReview}

% ── Warm-Up ──────────────────────────────────────────────────────────────────
\begin{practiceBox}[funGreen]{\faSun~Warm-Up}

\practiceHeader[funGreen]{Find the Actual Length}
Use the given scale to find each actual measurement.

\begin{multicols}{2}
\prob Scale: $1$ cm $= 3$ m. Drawing length: $5$ cm. \answerBlank[3cm]
\answerExplain{$15$ m}{Multiply drawing length by the scale factor: $5 \times 3 = 15$ m.}

\prob Scale: $1$ in $= 10$ ft. Drawing length: $7$ in. \answerBlank[3cm]
\answerExplain{$70$ ft}{Multiply drawing length by the scale factor: $7 \times 10 = 70$ ft.}

\prob Scale: $1$ cm $= 50$ km. Drawing length: $2$ cm. \answerBlank[3cm]
\answerExplain{$100$ km}{Multiply drawing length by the scale factor: $2 \times 50 = 100$ km.}

\prob Scale: $1$ in $= 4$ ft. Drawing length: $3.5$ in. \answerBlank[3cm]
\answerExplain{$14$ ft}{Multiply drawing length by the scale factor: $3.5 \times 4 = 14$ ft.}

\prob Scale: $1$ cm $= 25$ m. Drawing length: $6$ cm. \answerBlank[3cm]
\answerExplain{$150$ m}{Multiply drawing length by the scale factor: $6 \times 25 = 150$ m.}

\prob Scale: $1$ in $= 5$ mi. Drawing length: $8$ in. \answerBlank[3cm]
\answerExplain{$40$ mi}{Multiply drawing length by the scale factor: $8 \times 5 = 40$ mi.}
\end{multicols}
\end{practiceBox}

% ── Practice A: Finding the Scale Factor ─────────────────────────────────────
\begin{practiceBox}{Finding the Scale Factor}

\practiceHeader{Find the scale factor for each scale.}

\begin{multicols}{2}
\prob $1$ cm $= 20$ km \answerBlank[3cm]
\answerExplain{$20$}{The scale factor is the number of real units per drawing unit. Here, $1$ cm represents $20$ km, so the scale factor is $20$.}

\prob $1$ in $= 15$ ft \answerBlank[3cm]
\answerExplain{$15$}{Each inch on the drawing represents $15$ ft, so the scale factor is $15$.}

\prob $2$ cm $= 40$ m \answerBlank[3cm]
\answerExplain{$20$}{Divide to find the unit rate: $40 \div 2 = 20$ m per cm. The scale factor is $20$.}

\prob $3$ in $= 12$ ft \answerBlank[3cm]
\answerExplain{$4$}{Divide to find the unit rate: $12 \div 3 = 4$ ft per in. The scale factor is $4$.}

\prob $5$ cm $= 200$ km \answerBlank[3cm]
\answerExplain{$40$}{Divide to find the unit rate: $200 \div 5 = 40$ km per cm. The scale factor is $40$.}

\prob $4$ in $= 100$ mi \answerBlank[3cm]
\answerExplain{$25$}{Divide to find the unit rate: $100 \div 4 = 25$ mi per in. The scale factor is $25$.}
\end{multicols}
\end{practiceBox}

% ── Practice B: Working Backwards ────────────────────────────────────────────
\begin{practiceBox}[funTeal]{From Actual to Drawing}

\practiceHeader[funTeal]{Find the drawing length given the actual length and scale.}

\prob A hallway is $48$ ft long. Scale: $1$ in $= 8$ ft. Drawing length? \answerBlank[3cm]
\answerExplain{$6$ in}{Divide the actual length by the scale factor: $48 \div 8 = 6$ in.}

\prob A park is $300$ m wide. Scale: $1$ cm $= 50$ m. Drawing length? \answerBlank[3cm]
\answerExplain{$6$ cm}{Divide the actual width by the scale factor: $300 \div 50 = 6$ cm.}

\prob A road is $180$ km. Scale: $1$ cm $= 30$ km. Drawing length? \answerBlank[3cm]
\answerExplain{$6$ cm}{Divide the actual distance by the scale factor: $180 \div 30 = 6$ cm.}

\prob A building is $60$ ft tall. Scale: $1$ in $= 12$ ft. Drawing height? \answerBlank[3cm]
\answerExplain{$5$ in}{Divide the actual height by the scale factor: $60 \div 12 = 5$ in.}
\end{practiceBox}

% ── Practice C: Area in Scale Drawings ───────────────────────────────────────
\begin{practiceBox}[funPurple]{Scale and Area}

\practiceHeader[funPurple]{Remember: if lengths scale by $k$, areas scale by $k^2$.}

\prob A room on a blueprint is $3$ cm by $4$ cm. Scale: $1$ cm $= 5$ m. What is the actual area?
\answerBlank[3cm]
\answerExplain{$300$ m$^2$}{Actual dimensions: $3 \times 5 = 15$ m and $4 \times 5 = 20$ m. Actual area: $15 \times 20 = 300$ m$^2$.}

\prob A garden on a map is $2$ cm by $6$ cm. Scale: $1$ cm $= 10$ m. What is the actual area?
\answerBlank[3cm]
\answerExplain{$1{,}200$ m$^2$}{Actual dimensions: $2 \times 10 = 20$ m and $6 \times 10 = 60$ m. Actual area: $20 \times 60 = 1{,}200$ m$^2$.}

\prob A scale factor is $8$. A square on the drawing has area $9$ cm$^2$. What is the actual area?
\answerBlank[3cm]
\answerExplain{$576$ cm$^2$}{Area scales by $k^2 = 8^2 = 64$. Actual area: $9 \times 64 = 576$ cm$^2$.}

\prob A scale factor is $4$. The actual area of a rectangle is $320$ m$^2$. What is the drawing area?
\answerBlank[3cm]
\answerExplain{$20$ m$^2$}{Drawing area $= \text{actual area} \div k^2 = 320 \div 16 = 20$ m$^2$.}
\end{practiceBox}

% ── Practice D: True or False ────────────────────────────────────────────────
\begin{practiceBox}[funBlue]{True or False?}

\trueOrFalse{In a scale drawing, angles are changed along with lengths.}
\answerExplain{False}{A scale drawing changes all lengths by the same factor, but angles remain exactly the same. Only lengths and areas change.}

\trueOrFalse{If the scale factor is $10$, and a wall is $2$ cm on the drawing, the actual wall is $20$ m.}
\answerExplain{True}{Multiply drawing length by the scale factor: $2 \times 10 = 20$. The units depend on the scale, and here the answer is $20$ m.}

\trueOrFalse{If you double the scale factor, the area of every shape on the drawing quadruples.}
\answerExplain{False}{Doubling the scale factor doubles lengths in the real object. The relationship between scale factor and area means $k^2$ changes. If the original scale factor was $k$ and you double it to $2k$, the new area factor is $(2k)^2 = 4k^2$, which is $4$ times the old $k^2$, so actual areas quadruple — but the drawing area stays the same; it is the actual area interpretation that changes.}
\end{practiceBox}

% ── Challenge ────────────────────────────────────────────────────────────────
\begin{challengeBox}
\prob A map uses the scale $1$ cm $= 15$ km. Town A and Town B are $7.5$ cm apart on the map. A car drives at $90$ km/h. How long does the trip take?
\answerExplain{$1.25$ hours (or $1$ hour $15$ minutes)}{Actual distance: $7.5 \times 15 = 112.5$ km. Time $= \text{distance} \div \text{speed} = 112.5 \div 90 = 1.25$ hours, which is $1$ hour and $15$ minutes.}

\prob A rectangular pool on a blueprint is $5$ cm by $3$ cm. The scale is $1$ cm $= 4$ m. If the pool is $2$ m deep, what is the volume of the pool?
\answerExplain{$480$ m$^3$}{Actual dimensions: $5 \times 4 = 20$ m by $3 \times 4 = 12$ m. Volume $= 20 \times 12 \times 2 = 480$ m$^3$.}
\end{challengeBox}

\encouragement{Great work with scale drawings! You can now read blueprints and maps like a pro!}
